% Accepted after independent root, intervals, and schematization audits; integrated.
\begin{claim}{C.29}{Proved under stated hypotheses}{Finite Hom modules suffice for pointwise tensors}\label{C29}
Assume the normalized corollary hypotheses of Claim~\ref{C6}, including the full interval condition. Suppose there is a chosen presentation
\[
 \mathcal M\simeq\operatorname{Fun}_k(\mathcal A^{\mathrm{op}},\Mod k)
\]
by a small $k$-linear indexing category $\mathcal A$, not required to contain a zero object or finite sums, such that every $\mathcal A(a,b)$ is finitely generated projective over $k$. Suppose the given coherent internal tensor is pointwise over $k$, with unit value $\mathbf1(a)=k$ compatibly with the pointwise unit maps. This is a hypothesis on the tensor already supplied, not an assertion that tensoring two structure maps of arbitrary linear presheaves defines a $k$-linear presheaf. Then every chosen projective resolving adjunction $F\dashv G$ of Claim~\ref{C6} induces the desired equivalence. No finite-presentation assumption on $\mathbf1$ as an $\mathcal A$-module, no flatness assumption on $J$, and no countable scalar form is required.
\end{claim}
\begin{proof}
We construct a filtered presentation of $J$ in the category of objects under $\mathbf1$; the stages need not be projective or flat in $\mathcal M$. Choose a representable-generator epimorphism
\[
 \pi:\bigoplus_{t\in T}h_{a_t}\longrightarrow J.
\]
The map $[\alpha,\pi]:\mathbf1\oplus\bigoplus_{t\in T}h_{a_t}\to J$ is an epimorphism. Choose another sum of representables $\bigoplus_{r\in R}h_{b_r}$ mapping epimorphically onto its kernel. Write $d_r:h_{b_r}\to\mathbf1\oplus\bigoplus_{t\in T}h_{a_t}$ for the resulting relation map. Since a representable is finitely presented, $d_r$ factors through $\mathbf1\oplus\bigoplus_{t\in T_r}h_{a_t}$ for a finite subset $T_r\subseteq T$. This also follows directly by evaluating $d_r$ on the identity element: its component in the coproduct has finite support.

Let $I$ be the filtered poset of pairs $(T_0,R_0)$ of finite subsets of $T,R$ such that $T_r\subseteq T_0$ for every $r\in R_0$, ordered by inclusion. Set
\[
 P_{T_0,R_0}=\coker\left(
 \bigoplus_{r\in R_0}h_{b_r}\longrightarrow
 \mathbf1\oplus\bigoplus_{t\in T_0}h_{a_t}\right).
\]
Each stage carries a compatible map $\alpha_i:\mathbf1\to P_i$, and the presentation gives $\operatorname{colim}_I P_i=J$ under $\mathbf1$. Indeed the filtered colimits of the relation and generator diagrams are their indicated full coproducts, and colimits preserve the displayed cokernels. This argument does not require $\mathbf1$ itself to be finitely presented.

For every $X$, the composite of $\alpha_i\otimes X$ with $P_i\otimes X\to J\otimes X$ is $\alpha\otimes X$, hence $\alpha_i\otimes X$ is monic. Evaluating at $a$ and taking $X=\kappa(V)$ for arbitrary $V\in\Mod k$ shows that
\[
 \alpha_i(a):k\longrightarrow P_i(a)
\]
is pure. Moreover $P_i(a)$ is finitely presented over $k$: the displayed finite presentation of the stage evaluates to a cokernel between finite sums of finitely generated projective modules, with one additional summand $k$ in the target. Its quotient $P_i(a)/k$ is finitely presented as well, since we obtain a finite presentation by adding the single relation given by $\alpha_i(a)(1)$.

We recall explicitly why a pure inclusion with finitely presented quotient splits. In a pure exact sequence $0\to A\to B\to C\to0$ with $C$ finitely presented, choose finitely many generators of $C$, their lifts $y_j$ to $B$, and a finite matrix of relations. Applying a relation to the lifts gives constants in $A$. The resulting finite system of linear equations has the solution $y_j$ in $B$. Purity says that it has a solution in $A$. Subtract those elements of $A$ from the chosen lifts; the adjusted lifts satisfy the defining relations of $C$ and yield a section $C\to B$. Here the finite-equation characterization follows directly from tensor-purity: if a vector of constants in $A^m$ lies in the image of a finite matrix acting on $B^n$, its class in the tensor product with the finitely presented cokernel of that matrix is zero; injectivity of the tensor map says that it was already in the image of the matrix on $A^n$. This proves the assertion without a coherence hypothesis on $k$.

Thus each $\alpha_i(a)$ splits. Pointwise tensoring gives an evaluationwise splitting of $\alpha_i\otimes X$, and every component $F_n$, being a retract of a sum of evaluations, takes this map to a split monomorphism. No splitting is required to be natural in $a$ or in $i$.

Let $K$ be an acyclic nonnegative complex of projective $k$-modules and put $Q=GK$. Its disk decomposition makes it injective relative to degreewise split monomorphisms. Hence the counit $FQ\to K$ extends across $F(\alpha_i\otimes Q)$ for each $i$. Adjunction gives a retraction
\[
 r_i:P_i\otimes Q\longrightarrow Q,
 \qquad r_i(\alpha_i\otimes Q)=\id_Q.
\]
Pointwise tensoring and $F$ preserve filtered colimits, so
\[
 \operatorname{colim}_i F(P_i\otimes Q)
 \cong F(J\otimes Q)\simeq0.
\]
Homology commutes with filtered colimits of $k$-modules. A class $z\in H_n(FQ)$ maps to some stage and then vanishes at a later stage. Compatibility of the $\alpha_i$ identifies its later image with $H_nF(\alpha_j\otimes Q)(z)$. Applying $H_nF(r_j)$ gives $z=0$. As in Claim~\ref{C27}, compatibility of the retractions and nullity of individual stages are unnecessary.

We have proved the acyclic-projective test of Claim~\ref{C10}. The canonical normalization, coproducts, and retracts prove the degree-zero projective-coefficient test exactly as in Claim~\ref{C7}. Claim~\ref{C3} now proves the raw conclusion for every chosen resolving adjunction.
\end{proof}
